\documentclass[a4paper,12pt]{article}

\usepackage{iftex}
\ifPDFTeX
  \usepackage[T2A]{fontenc}
  \usepackage[utf8]{inputenc}
  \usepackage[english,russian]{babel}
\else
  \usepackage{fontspec}
  \setmainfont{DejaVu Serif}
  \setsansfont{DejaVu Sans}
  \setmonofont{DejaVu Sans Mono}[Scale=0.9]
  \usepackage{polyglossia}
  \setdefaultlanguage{russian}
  \setotherlanguage{english}
\fi

\usepackage[left=25mm,right=15mm,top=20mm,bottom=20mm]{geometry}
\usepackage{indentfirst}
\usepackage{setspace}
\onehalfspacing
\emergencystretch=3em
\usepackage{amsmath,amssymb}
\usepackage{array,longtable,booktabs,tabularx}
\usepackage{enumitem}
\setlist{nosep,leftmargin=*}
\usepackage[table]{xcolor}
\usepackage{listings}
\usepackage{tikz}
\usetikzlibrary{shapes.geometric,positioning,arrows.meta,fit,calc}
\usepackage{float}
\usepackage{adjustbox}
\usepackage{caption}
\captionsetup{labelsep=endash}
\usepackage[hidelinks]{hyperref}

\lstset{
  basicstyle=\ttfamily\footnotesize,
  backgroundcolor=\color[RGB]{246,246,246},
  frame=single, rulecolor=\color{black!20},
  breaklines=true, columns=fullflexible,
  keepspaces=true, showstringspaces=false,
  extendedchars=true
}

% --- элементы диаграмм прецедентов ---
\tikzset{
  uc/.style={ellipse, draw, align=center, minimum width=50mm,
             minimum height=12mm, font=\small, inner sep=1pt,
             execute at begin node=\setstretch{1}},
  sys/.style={draw, rounded corners, inner sep=5mm},
  assoc/.style={-, thin},
  rel/.style={-{Stealth[length=2mm]}, dashed, thin},
  gen/.style={-{Triangle[open,length=3mm]}, thin},
  lbl/.style={font=\scriptsize, fill=white, inner sep=1pt}
}
% актор: \actor{имя узла}{позиция}{подпись}
\newcommand{\actor}[3]{%
  \begin{scope}[shift={#2}]
    \draw (0,0.55) circle (0.18);
    \draw (0,0.37) -- (0,-0.05);
    \draw (-0.3,0.22) -- (0.3,0.22);
    \draw (0,-0.05) -- (-0.22,-0.4);
    \draw (0,-0.05) -- (0.22,-0.4);
    \node[font=\small, align=center, below] at (0,-0.42) {#3};
    \coordinate (#1) at (0,0.15);
    \coordinate (#1-r) at (0.4,0.2);
    \coordinate (#1-l) at (-0.4,0.2);
  \end{scope}}

% Требование: \req{ID}{текст}
\newcommand{\req}[2]{\textbf{#1} & #2 \\ \hline}
\newenvironment{reqtable}{%
  \begin{longtable}{|>{\raggedright\arraybackslash}p{19mm}|>{\raggedright\arraybackslash}p{137mm}|}
  \hline
  \rowcolor{black!8}\textbf{ID} & \textbf{Требование} \\ \hline
  \endhead}{\end{longtable}}

% Таблица описания прецедента
\newcommand{\ucrow}[2]{\textbf{#1} & #2 \\ \hline}
\newenvironment{uctable}{%
  \begin{longtable}{|>{\raggedright\arraybackslash}p{38mm}|>{\raggedright\arraybackslash}p{118mm}|}
  \hline}{\end{longtable}\vspace{-4mm}\addtocounter{table}{-1}}
\newcommand{\steps}[1]{\begin{minipage}[t]{\linewidth}\begin{enumerate}[nosep,leftmargin=*]#1\end{enumerate}\vspace{1mm}\end{minipage}}
\tikzset{cond/.style={font=\scriptsize, fill=white, inner sep=1pt, align=center}}

% Описание прецедента
\newenvironment{usecase}[2]{%
  \subsubsection*{#1. #2}
  \begin{description}[leftmargin=0pt,labelindent=0pt,style=nextline]}
  {\end{description}}

\author{}
\date{}
\begin{document}


\begin{titlepage}

\thispagestyle{empty}

\centerline{Университет ИТМО}
\centerline{Факультет программной инженерии и компьютерной техники}


\vfill

\centerline{\huge{Отчет}}
\centerline{\large{по курсовой работе}}
\centerline{\LARGE{Информационная система «Хогвартс»}}
\centerline{\large{Этап №1}}

\vfill

Группа  \hfill P3309

Выполнили: \hfill Исаева Александра-Ирина Антоновна\hfill


Преподаватель \hfill  Николаев Владимир Вячеславович

\vfill

\centerline{Санкт-Петербург, 2026}
\clearpage
\end{titlepage}
\newpage
\setcounter{page}{2}
 
\tableofcontents
\newpage
\section{Описание предметной области}


\subsection{Общие сведения}
 
Хогвартс — школа-интернат для детей от 11 до 17 лет, обучение длится семь курсов.
Каждый ученик при поступлении распределяется на один из четырёх факультетов
(Гриффиндор, Слизерин, Когтевран, Пуффендуй) и остаётся на нём до выпуска.
У факультета есть декан из числа преподавателей, гостиная и спальни, разделённые
по курсу и полу. Занятия проходят по расписанию, составленному для курса и факультета.
 
Факультеты соревнуются за Кубок школы: преподаватели начисляют и снимают очки
факультету за поступки его учеников. За порядком следит смотритель Аргус Филч:
он фиксирует нарушения, принимает отработки, ведёт совятню, досматривает посылки
и проверяет учеников при выходе из замка. Здоровьем учеников занимается
целитель мадам Помфри; лечебные зелья для больничного крыла варит преподаватель
зельеварения из ингредиентов, выращенных в теплицах травологии или купленных в лавках
Хогсмида. Связь с внешним миром идёт через совятню: письма и посылки доставляют совы.
С 3 курса ученики с разрешения опекуна могут ходить в соседнюю деревню Хогсмид.
 
Пользователями системы являются сотрудники школы. Ученики в системе — объекты учёта:
их распределяют, заселяют, наказывают, лечат и допускают в Хогсмид сотрудники,
а о принятых решениях ученикам сообщают декан, преподаватель или смотритель.
 
Основные объекты предметной области: ученик, факультет, спальня, занятие, очки
факультета, нарушение, отработка, сова, почтовое отправление, конфискованный предмет,
обращение в больничное крыло, койка, зелье, ингредиент, заявка на варку зелья,
выход в Хогсмид, разрешение опекуна.
 
\subsection{Роли пользователей}
 
Работа с системой начинается со входа: неавторизованный пользователь может только
ввести логин и пароль. После входа система определяет роль пользователя. 
 
\begin{longtable}{|>{\raggedright\arraybackslash}p{42mm}|>{\raggedright\arraybackslash}p{112mm}|}
\hline
\rowcolor{black!8}\textbf{Роль} & \textbf{Что делает в системе} \\ \hline
\endhead
Директор & Утверждение списка первокурсников и формирование писем о зачислении, ввод решения Распределяющей шляпы, расписание занятий, объявление выходов в Хогсмид, подведение итогов года и вручение Кубка школы, отчёты, учётные записи и справочники \\ \hline
Декан & Преподаватель, отвечающий за факультет: всё, что доступно преподавателю, а для своего факультета — заселение учеников, утверждение разрешений на Хогсмид и запрет выхода, контроль нарушений и отработок, снятие очков по нарушениям, зафиксированным смотрителем \\ \hline
Преподаватель & Фиксация нарушений и назначение отработок по ним, начисление и снятие очков, просмотр освобождений учеников от занятий \\ \hline
Смотритель & Фиксация нарушений и назначение отработок по ним, приём отработок, учёт сов, отправка писем от имени школы, приём и досмотр посылок, реестр конфиската, отметка выхода и возвращения из Хогсмида \\ \hline
Целитель & Журнал обращений, койки, запас зелий, заявки на варку \\ \hline
Зельевар & Варка зелий по заявкам, учёт ингредиентов (урожай теплиц и поставки), заказы поставщикам Хогсмида \\ \hline
\end{longtable}
 
Ученики, Распределяющая шляпа, родители и опекуны и поставщики Хогсмида
пользователями системы не являются: решения по ним и полученные от них документы
вносят в систему сотрудники школы.
 
\subsection{Бизнес-процессы Хогвартса}
 
 
% Жизнь Хогвартса держится на нескольких сквозных процессах. В каждом участвует
% несколько ролей, а одно событие вызывает последствия сразу в нескольких подсистемах.
% Эти процессы характерны именно для Хогвартса и в обычной школе не встречаются.
 
\subsubsection*{бп1. Распределение первокурсников}
 
\begin{description}[leftmargin=0pt,labelindent=0pt,style=nextline]
\item[Участники] Директор (утверждает список первокурсников, формирует письма о зачислении и проводит церемонию),
  смотритель (отправляет письма о зачислении совами), Распределяющая шляпа (выбирает
  факультет), декан факультета (заселяет ученика);
  распределяемый первокурсник.
\item[Ход процесса]\leavevmode
  \begin{enumerate}
        \item Директор утверждает список будущих первокурсников и формирует по нему письма о зачислении; смотритель отправляет их совиной почтой.
    \item Первокурсники прибывают в школу на Хогвартс-экспрессе.
    \item На церемонии в Большом зале Распределяющая шляпа выбирает факультет с учётом пожеланий ученика.
    \item Декан факультета заселяет ученика в спальню своего факультета, курса и пола.
    \item Ученику назначается расписание его курса и факультета; с этого момента его поступки приносят или отнимают очки факультету.
  \end{enumerate}
 
\item[Роль ИС] Формирует письма о зачислении по утверждённому списку и передаёт их на отправку, хранит решение о распределении, подбирает свободное место в спальне
  нужного факультета, пола и курса, сразу привязывает ученика к расписанию и рейтингу
  факультета.
\end{description}
 
\subsubsection*{бп2. Дисциплина}
 
\begin{description}[leftmargin=0pt,labelindent=0pt,style=nextline]
\item[Участники] Преподаватель или смотритель Аргус Филч (фиксирует нарушение
  и назначает отработку), смотритель (принимает отработку), декан факультета
  нарушителя (контролирует дисциплину); ученик-нарушитель.
\item[Ход процесса]\leavevmode
  \begin{enumerate}
    \item Преподаватель или смотритель застаёт ученика на нарушении: в коридоре после отбоя, в Запретном лесу, с запрещённым предметом.
    \item Преподаватель снимает очки с факультета ученика. По нарушению, которое зафиксировал смотритель, очки снимает декан.
    \item Тот, кто зафиксировал нарушение, исходя из его причины назначает отработку: чистку кубков в Зале трофеев без магии, помощь леснику в Запретном лесу, разбор архива.
    \item В назначенное время Филч принимает отработку и отмечает её выполнение.
    \item Если ученик не явился, фиксируется новое нарушение и цикл повторяется.
  \end{enumerate}
 
\item[Роль ИС] Связывает нарушение, списание очков и отработку в одну цепочку,
  проверяет права ролей, не даёт назначить отработку на время урока,
  ведёт рейтинг факультетов с полной историей изменений.
\end{description}
 
\subsubsection*{бп3. Совиная почта с досмотром}
 
\begin{description}[leftmargin=0pt,labelindent=0pt,style=nextline]
\item[Участники] Смотритель Филч (ведёт совятню: отправляет письма от имени школы,
  принимает и досматривает посылки), отправители посылок (родители, магазины Хогсмида
  «Зонко» и «Сладкое королевство»); ученик-получатель.
\item[Ход процесса]\leavevmode
  \begin{enumerate}
    \item Посылка ученику прилетает совой в совятню, Филч регистрирует её: отправитель, получатель, вес, сова.
        \item Филч отправляет письма от имени школы (например, письма о зачислении первокурсникам): для каждого письма подбирается свободная сова, способная его унести.
    \item Посылка без нарушений выдаётся ученику.
    \item Запрещённый предмет конфискуется и заносится в реестр конфиската, на ученика оформляется нарушение, и запускается бп2.
    \item В конце учебного года конфискат возвращают владельцу или уничтожают.
  \end{enumerate}
\item[Особенность] Почту доставляют живые совы с ограничением по весу
  и времени отдыха, существует отдельный перечень запрещённых магических предметов,
  а досмотр посылок — обязанность смотрителя.
\item[Роль ИС] Подбирает свободную сову по весу письма или посылки, отслеживает статус
  доставки, ведёт реестр конфиската и автоматически передаёт нарушение
  в дисциплинарный цикл.
\end{description}
 
\subsubsection*{бп4. Лечение и обеспечение больничного крыла зельями}
 
\begin{description}[leftmargin=0pt,labelindent=0pt,style=nextline]
\item[Участники] Целитель мадам Помфри (лечит), зельевар — преподаватель
  зельеварения (варит лечебные зелья, учитывает урожай ингредиентов из теплиц
  травологии), поставщики Хогсмида (недостающие ингредиенты), преподаватели
  (видят освобождение ученика); пострадавший ученик.
\item[Ход процесса]\leavevmode
  \begin{enumerate}
    \item Ученик получает травму на матче по квиддичу, уроке ухода за магическими существами, на дуэли или заболевает и попадает в больничное крыло.
    \item Целитель принимает ученика, при необходимости кладёт на койку и лечит зельями (Костерост, Бодроперцовое зелье). Запас зелий уменьшается.
    \item На время госпитализации ученик освобождён от занятий, преподаватели видят это и не ставят ему оценки.
    \item Когда запас зелья опускается ниже минимума, целитель отправляет заявку преподавателю зельеварения.
    \item Зельевар варит зелье из ингредиентов, выращенных в теплицах травологии; недостающие он заказывает у поставщиков Хогсмида.
    \item Готовое зелье поступает в больничное крыло, запас пополняется.
  \end{enumerate}
\item[Особенность] Лекарства не покупают, а варит преподаватель школы,
  ингредиенты растут в школьных теплицах, а травмы магические: квиддич, существа, заклятия.
\item[Роль ИС] Ведёт журнал травм по причинам, учитывает свободные койки,
  автоматически освобождает ученика от занятий, следит за запасом зелий и формирует
  заявку зельевару и заказ поставщику.
\end{description}
 
\subsubsection*{бп5. Выход в Хогсмид}
 
\begin{description}[leftmargin=0pt,labelindent=0pt,style=nextline]
\item[Участники] Директор (объявляет выход), родитель или опекун (подписывает
  разрешение и присылает его совой), декан (утверждает разрешение), смотритель Филч
  (проверяет учеников по списку на выходе и при возвращении); ученик 3 курса и старше.
\item[Ход процесса]\leavevmode
  \begin{enumerate}
    \item Директор объявляет дату выхода в Хогсмид.
    \item Родитель или опекун присылает подписанное разрешение совиной почтой, оно передаётся декану.
    \item Декан утверждает разрешение. Ученику с невыполненной отработкой выход может быть запрещён.
    \item В день выхода формируется список допущенных учеников, Филч проверяет учеников по списку.
    \item При возвращении Филч отмечает время; опоздание фиксируется как нарушение.
  \end{enumerate}
\item[Особенность] Выход связывает три подсистемы: почту (наличие разрешения), дисциплину (запрет за нарушения) и смотрителя на выходе.
  Допуск зависит от курса, подписи опекуна и поведения ученика.
\item[Роль ИС] Хранит разрешения, проверяет курс ученика и наличие невыполненных
  отработок, формирует список допущенных для Филча и фиксирует опоздания как нарушения.
\end{description}
 
% \subsubsection*{Связи между процессами}
 
% \begin{itemize}
%   \item бп1 $\to$ все остальные: факультет, выбранный при распределении, определяет, чьи очки меняются в бп2 и кто декан ученика в бп4 и бп5;
%   \item бп3 $\to$ бп2: запрещённый предмет в посылке приводит к нарушению;
%   \item бп3 $\to$ бп5: разрешение на выход приходит совиной почтой;
%   \item бп2 $\to$ бп5: невыполненная отработка закрывает выход в Хогсмид;
%   \item бп4 $\to$ учебный процесс: госпитализированный ученик освобождён от занятий.
% \end{itemize}
 
% =====================================================================
\section{Назначение информационной системы}
% =====================================================================
 
% Сейчас каждый шаг этих процессов выполняется вручную: нарушения записываются на
% пергаменте, очки считаются по песочным часам факультетов, разрешения на Хогсмид
% хранятся у декана, запас зелий знает только мадам Помфри. Поэтому цепочки рвутся:
% отработку назначают на время урока, пропущенная отработка остаётся без последствий,
% конфискованный предмет не связан с нарушением, а целитель узнаёт о нехватке зелья,
% когда оно уже закончилось.
 
\textbf{Цель ИС} — автоматизировать сквозные процессы жизни Хогвартса, в которых
участвуют несколько ролей: распределение первокурсников, дисциплинарный цикл,
доставку почты совами с досмотром посылок, лечение учеников с обеспечением
больничного крыла зельями и выход в Хогсмид. Каждое событие должно сразу
отражаться во всех связанных подсистемах — очках факультетов, расписании учеников,
запасах зелий и списках допуска, — а каждая роль выполнять только свой шаг процесса.
 
\textbf{Задачи, которые решает система:}
\begin{enumerate}
  \item Формировать письма о зачислении по утверждённому списку первокурсников,
        фиксировать решение Распределяющей шляпы и заселять первокурсника
        в спальню своего факультета, пола и курса с проверкой свободных мест.
  \item Фиксировать нарушение и в той же операции списывать очки с факультета
        нарушителя с учётом прав роли.
  \item Назначать отработку с проверкой занятости ученика, контролировать её
        выполнение и при пропуске автоматически создавать новое нарушение.
  \item Вести рейтинг факультетов в борьбе за Кубок школы с полной историей
        изменений: кто, за что и сколько очков.
  \item Регистрировать отправления, подбирать сову по весу посылки и её
        состоянию и отслеживать статус доставки.
  \item Организовать досмотр посылок, вести реестр конфиската и при обнаружении
        запрещённого предмета запускать дисциплинарный цикл.
  \item Регистрировать травмы с причиной, учитывать койки и автоматически
        освобождать ученика от занятий.
  \item Учитывать запас лечебных зелий, формировать заявку на варку зельевару
        и заказ недостающих ингредиентов поставщикам Хогсмида.
  \item Хранить разрешения опекунов, проверять допуск в Хогсмид по курсу
        и невыполненным отработкам, формировать список для смотрителя
        и фиксировать опоздания.
  \item Разграничивать права так, чтобы каждая роль выполняла только свой шаг
        процесса.
  \item Формировать для директора и деканов отчёты по каждому процессу:
        нарушения, травматизм, расход зелий, итоги Кубка школы.
\end{enumerate}
 
\section{Требования к системе}
 
\subsection{Функциональные требования}
 
 
\subsubsection{Общие функции}
\begin{reqtable}
\req{FR-01}{Система должна позволять неавторизованному пользователю выполнить вход по логину и паролю.}
\req{FR-02}{Система должна отклонять вход при неверном логине или пароле, а также для заблокированной учётной записи.}
\req{FR-03}{Система должна позволять авторизованному пользователю завершить сеанс.}
\req{FR-04}{После входа система должна определять роль пользователя и предоставлять ему доступ только к функциям этой роли.}
\req{FR-05}{Система должна позволять директору создавать и блокировать учётные записи сотрудников и назначать им роли.}
\req{FR-06}{Система должна позволять директору вести справочники: факультеты, спальни, типы нарушений, запрещённые предметы, зелья и их рецепты, ингредиенты.}
\req{FR-07}{Система должна хранить сведения об учениках (ФИО, дата рождения, пол, курс, факультет), сотрудниках, спальнях, расписании занятий, совах и поставщиках Хогсмида.}
\req{FR-08}{Система должна позволять директору вести расписание занятий по курсам и факультетам: предмет, преподаватель, аудитория, день, время.}
\req{FR-09}{Система должна показывать каждому пользователю главную страницу со сводкой по его роли: декану — нарушения и отработки своего факультета; смотрителю — отработки на сегодня и посылки на досмотр; целителю — занятые койки и зелья ниже минимума; зельевару — открытые заявки на варку.}
\req{FR-10}{Система должна предоставлять каждой роли доступ к любой её функции не более чем за 3 перехода с главной страницы роли.}
\req{FR-11}{Перед необратимыми действиями (изъятие предмета, закрытие учебного года, блокировка учётной записи) система должна запрашивать подтверждение пользователя.}
\end{reqtable}
 
\subsubsection{Распределение первокурсников}
\begin{reqtable}
\req{FR-12}{Система должна позволять директору формировать и утверждать список будущих первокурсников. По утверждённому списку система формирует письма о зачислении и передаёт их в очередь отправки смотрителя.}
\req{FR-13}{Система должна позволять директору вносить решение Распределяющей шляпы: ученик, факультет, дата.}
\req{FR-14}{Система не должна допускать повторного распределения ученика.}
\req{FR-15}{После распределения система должна автоматически привязывать ученика к расписанию его курса и факультета.}
\req{FR-16}{Система должна предоставлять декану список спален своего факультета, подходящих ученику по полу и курсу и имеющих свободные места.}
\req{FR-17}{Система должна позволять декану заселить ученика в выбранную спальню.}
\req{FR-18}{Система не должна допускать заселения в спальню больше учеников, чем в ней мест, в том числе когда несколько пользователей заселяют учеников одновременно.}
\end{reqtable}
 
\subsubsection{Дисциплинарный цикл}
\begin{reqtable}
\req{FR-19}{Система должна позволять преподавателю (в том числе декану) и смотрителю регистрировать нарушение: ученик, тип нарушения, место, время, описание.}
\req{FR-20}{Система должна позволять преподавателю (в том числе декану) и директору начислять и снимать очки факультету с указанием ученика, суммы и причины.}
\req{FR-21}{Система должна позволять преподавателю при регистрации нарушения сразу снять очки с факультета нарушителя. Нарушение и снятие очков должны сохраняться вместе: если одно из них не сохранилось, не сохраняется ничего.}
\req{FR-22}{Система должна позволять декану снять очки по нарушению ученика своего факультета, которое зарегистрировал смотритель или система.}
\req{FR-23}{Система должна позволять пользователю, регистрирующему нарушение, в той же операции назначить по нему отработку: дата, время, место, вид работы. Назначить отработку без зарегистрированного нарушения нельзя.}
\req{FR-24}{Система не должна допускать назначения отработки на время урока ученика или другой его отработки.}
\req{FR-25}{Система должна позволять смотрителю отмечать выполнение отработки, её перенос или неявку ученика.}
\req{FR-26}{При неявке на отработку система должна автоматически регистрировать новое нарушение.}
\req{FR-27}{Система должна показывать рейтинг факультетов с учётом всех сохранённых изменений очков.}
\req{FR-28}{Система должна хранить журнал изменений очков (кто изменил, у какого ученика, сколько, причина, время) и не допускать изменения или удаления его записей.}
\req{FR-29}{Система должна позволять фильтровать журнал очков по факультету, ученику и автору изменения.}
\req{FR-30}{Система должна показывать декану нарушения и отработки учеников своего факультета.}
\req{FR-31}{Система должна позволять директору закрыть учебный год и зафиксировать победителя Кубка школы.}
\end{reqtable}
 
\subsubsection{Совиная почта с досмотром}
\begin{reqtable}
\req{FR-32}{Система должна позволять смотрителю вести учёт сов: школьная или личная, владелец, максимальный вес груза, статус (свободна, в полёте, отдыхает).}
\req{FR-33}{Система должна позволять смотрителю регистрировать входящие посылки и исходящие письма от имени школы (в том числе письма о зачислении первокурсникам): отправитель, получатель, тип, вес, сова.}
\req{FR-34}{Для исходящего отправления система должна предлагать только свободных сов, способных унести его вес.}
\req{FR-35}{Если подходящей совы нет, система должна ставить отправление в очередь.}
\req{FR-36}{Система должна позволять смотрителю отмечать доставку или возврат отправления. После доставки сова переводится в статус «отдыхает».}
\req{FR-37}{Система должна помещать входящие посылки ученикам в очередь досмотра и переводить посылку в статус «допущена к выдаче» только после досмотра.}
\req{FR-38}{Система должна позволять смотрителю отметить посылку как проверенную или изъять из неё предметы, внесённые в справочник запрещённых.}
\req{FR-39}{При изъятии предмета система должна заносить его в реестр конфиската и автоматически регистрировать нарушение на получателя посылки.}
\req{FR-40}{Система должна позволять смотрителю в конце учебного года отметить возврат или уничтожение конфиската.}
\end{reqtable}
 
\subsubsection{Лечение и обеспечение зельями}
\begin{reqtable}
\req{FR-41}{Система должна позволять целителю регистрировать обращение: пациент, травма, причина (квиддич, урок, дуэль, магические существа, прочее), назначенное лечение.}
\req{FR-42}{Система должна позволять целителю госпитализировать ученика на свободную койку и выписывать его.}
\req{FR-43}{Система не должна допускать госпитализации двух учеников на одну койку, в том числе при одновременных действиях.}
\req{FR-44}{На время госпитализации система должна отмечать ученика освобождённым от занятий и показывать эту отметку преподавателям.}
\req{FR-45}{Медицинские данные ученика система должна показывать только целителю, декану факультета ученика и директору; преподаватели видят только факт освобождения от занятий.}
\req{FR-46}{Система должна позволять целителю списывать выданные зелья и не допускать отрицательного запаса.}
\req{FR-47}{Когда запас зелья опускается ниже минимума, система должна предлагать целителю создать заявку на варку.}
\req{FR-48}{Система должна позволять целителю создавать заявку на варку: зелье, количество, желаемый срок.}
\req{FR-49}{Система должна показывать зельевару открытые заявки и потребность в ингредиентах по рецепту.}
\req{FR-50}{Система должна позволять зельевару отметить варку зелья. При этом ингредиенты списываются по рецепту, запас зелья пополняется и заявка закрывается одной операцией.}
\req{FR-51}{Если ингредиентов не хватает, система должна показывать дефицит и переводить заявку в статус «ожидает ингредиентов».}
\req{FR-52}{Система должна позволять зельевару регистрировать урожай ингредиентов из теплиц. Урожай поступает в общий запас ингредиентов, из которого списываются ингредиенты при варке зелий.}
\req{FR-53}{Система должна позволять зельевару оформлять заказ ингредиентов поставщику Хогсмида и регистрировать поставку с фактическим количеством.}
\req{FR-54}{Система должна показывать директору статистику травм по причинам.}
\end{reqtable}
 
\subsubsection{Выход в Хогсмид}
\begin{reqtable}
\req{FR-55}{Система должна позволять директору объявлять выход в Хогсмид: дата, время выхода, время возвращения.}
\req{FR-56}{Система должна позволять декану регистрировать полученное от опекуна разрешение и утверждать его. Разрешение действует один учебный год.}
\req{FR-57}{Система должна автоматически определять допуск ученика: 3 курс или старше, утверждённое разрешение, нет невыполненных отработок.}
\req{FR-58}{Система должна позволять декану запретить выход конкретному ученику с указанием причины.}
\req{FR-59}{Система должна формировать для смотрителя список допущенных учеников.}
\req{FR-60}{Система должна позволять смотрителю отмечать выход и возвращение каждого ученика и не допускать отметки выхода ученика, которого нет в списке.}
\req{FR-61}{Если ученик вернулся позже назначенного времени, система должна автоматически регистрировать нарушение.}
\end{reqtable}
 
\subsubsection{Отчётность}
\begin{reqtable}
\req{FR-62}{Система должна формировать для директора и деканов отчёты: нарушения и отработки по факультетам, травматизм по причинам, расход зелий и ингредиентов, итоги Кубка школы.}
\end{reqtable}
 
\subsection{Нефункциональные требования}
 
\subsubsection{Безопасность}
\begin{reqtable}
\req{NFR-01}{Пароли пользователей должны храниться только в виде стойких хешей с индивидуальной солью, без возможности восстановить исходный пароль.}
\req{NFR-02}{Ограничения доступа, установленные для ролей, должны сохраняться при прямом обращении к серверу, изменении адреса запроса и подмене идентификаторов объектов.}
\req{NFR-03}{Ввод SQL-кода или скриптов в поля форм не должен приводить к выполнению команд в базе данных или программного кода в браузере при последующем отображении.}
\req{NFR-04}{Запросы на изменение данных, отправленные со сторонних сайтов от имени вошедшего пользователя, должны отклоняться.}
\req{NFR-05}{Пароли, идентификаторы сеансов и учётные данные подключения к базе данных не должны попадать в журналы приложения.}
\end{reqtable}
 
\subsubsection{Надёжность и целостность данных}
\begin{reqtable}
\req{NFR-06}{Данные, успешно сохранённые пользователем, не должны теряться после завершения сеанса и после перезапуска приложения.}
\req{NFR-07}{После перезапуска приложения работоспособность должна восстанавливаться не более чем за 5 минут.}
\req{NFR-08}{При сбое во время выполнения операции в базе данных не должно оставаться частично сохранённых изменений.}
\end{reqtable}
 
\subsubsection{Производительность}
\begin{reqtable}
\req{NFR-09}{Страница со списком из 50 записей (нарушения, отправления, обращения) должна формироваться не более чем за 2 секунды.}
\req{NFR-10}{Сохранение нарушения со снятием очков вместе с обновлением рейтинга факультетов должно занимать не более 1 секунды.}
\req{NFR-11}{Формирование списка допуска на выход в Хогсмид для всех учеников 3--7 курсов должно занимать не более 3 секунд.}
\req{NFR-12}{Время отклика из NFR-09--NFR-11 должно выдерживаться при 50 одновременно работающих пользователях.}
\end{reqtable}
 
\subsubsection{Удобство использования}
\begin{reqtable}
\req{NFR-13}{Интерфейс должен быть на русском языке и иметь единообразную структуру навигации для всех ролей.}
\req{NFR-14}{Однотипные элементы управления должны иметь одинаковый вид и поведение.}
\req{NFR-15}{Сообщения об отказе должны объяснять причину понятным пользователю языком.}
\end{reqtable}
 
\subsubsection{Совместимость и условия эксплуатации}
\begin{reqtable}
\req{NFR-16}{Интерфейс должен корректно отображаться в последних версиях Chrome и Firefox без установки расширений.}
\req{NFR-17}{Приложение должно потреблять не более 512 МБ оперативной памяти и использовать не более 5 одновременных соединений с базой данных.}
\end{reqtable}
 
\section{Модели основных прецедентов}
 
Выделено шесть основных прецедентов: авторизация и по одному прецеденту на каждый
бизнес-процесс. Декан — частный случай актора «Преподаватель» (обобщение акторов):
ему доступны все прецеденты преподавателя.
 
 
\subsection{Модель прецедентов}
\begin{figure}[H]
\centering
\adjustbox{max width=\linewidth}{%
\begin{tikzpicture}
  \node[uc] (u1) at (0,10)  {Авторизация\\пользователя};
  \node[uc] (u2) at (0,8)   {Распределение и\\ заселение первокурсника};
  \node[uc] (u3) at (0,6)   {Нарушение\\и отработка};
  \node[uc] (u4) at (0,4)   {Совиная почта\\и досмотр посылок};
  \node[uc] (u5) at (0,2)   {Лечение и пополнение\\запаса зелий};
  \node[uc] (u6) at (0,0)   {Допуск и выход\\в Хогсмид};
  \node[sys, fit=(u1)(u6), label={[font=\small\bfseries]above:ИС «Хогвартс»}] {};

  \actor{anon}{(-8,10.0)}{Неавторизованный\\пользователь}
  \actor{dir}{(-8,7.2)}{Директор}
  \actor{dean}{(-8,4.0)}{Декан}
  \actor{prof}{(-8,0.8)}{Преподаватель}
  \actor{filch}{(8,7.6)}{Смотритель}
  \actor{heal}{(8,3.4)}{Целитель}
  \actor{pot}{(8,0.6)}{Зельевар}

  \draw[assoc] (anon-r) -- (u1.west);
  \draw[assoc] (dir-r) -- (u2.west);
  \draw[assoc] (dean-r) -- (u2.west);
  \draw[assoc] (dean-r) -- (u3.west);
  \draw[assoc] (dean-r) -- (u6.west);
  \draw[assoc] (prof-r) -- (u3.west);
  \draw[gen] (-8,3.0) -- (-8,1.6);
  \draw[assoc] (filch-l) -- (u3.east);
  \draw[assoc] (filch-l) -- (u4.east);
  \draw[assoc] (filch-l) -- (u6.east);
  \draw[assoc] (heal-l) -- (u5.east);
  \draw[assoc] (pot-l) -- (u5.east);
\end{tikzpicture}}
\caption{Модель прецедентов}
\end{figure}
 
Неавторизованный пользователь — не роль, а пользователь до входа: ему доступна
только авторизация. Все роли работают с системой после выполнения UC-01.
 
\setcounter{table}{0}
\subsection{Авторизация пользователя}
 
\begin{uctable}
\ucrow{ID}{UC-01}
\ucrow{Название}{Авторизация пользователя}
\ucrow{Краткое описание}{Пользователь проходит проверку учётных данных и получает доступ к функциям своей роли.}
\ucrow{Главный актор}{Неавторизованный пользователь}
\ucrow{Второстепенные акторы}{Нет}
\ucrow{Предусловия}{Учётная запись создана и не заблокирована.}
\ucrow{Основной поток}{\steps{
  \item Пользователь открывает страницу входа.
  \item Пользователь вводит логин и пароль.
  \item Система проверяет учётные данные и определяет роль пользователя.
  \item Система открывает главную страницу роли.}}
\ucrow{Альтернативные потоки}{Логин или пароль неверны либо учётная запись заблокирована — система выводит сообщение и предлагает повторить ввод.}
\ucrow{Постусловия}{Пользователь авторизован и работает в рабочей области своей роли.}
\ucrow{Точки расширения}{Отказ во входе}
\end{uctable}
\captionof{table}{Описание прецедента UC-01}
 
\begin{figure}[H]
\centering
\adjustbox{max width=0.85\linewidth}{%
\begin{tikzpicture}
  \node[uc] (a) at (0,3.2)   {Авторизоваться};
  \node[uc] (b) at (9.0,4.4) {Проверить\\учётные данные};
  \node[uc] (c) at (9.0,2.0) {Отказать во входе};
  \node[uc] (d) at (0,0)     {Завершить сеанс};
  \node[sys, fit=(a)(b)(c)(d), label={[font=\small\bfseries]above:UC-01. Авторизация пользователя}] {};
  \actor{anon}{(-6,3.0)}{Неавторизованный\\пользователь}
  \actor{auth}{(-6,-0.2)}{Авторизованный\\пользователь}
  \draw[assoc] (anon-r) -- (a.west);
  \draw[assoc] (auth-r) -- (d.west);
  \draw[rel] (a) -- node[lbl,above,sloped]{«include»} (b);
  \draw[rel] (c) -- node[cond,below,sloped]{«extend»\\{}[неверные данные]} (a);
\end{tikzpicture}}
\caption{Диаграмма прецедента UC-01}
\end{figure}
 
% ---------------------------------------------------------------------
\subsection{Распределение и заселение первокурсника}
 
\begin{uctable}
\ucrow{ID}{UC-02}
\ucrow{Название}{Распределение и заселение первокурсника}
\ucrow{Краткое описание}{Директор утверждает список первокурсников, система формирует письма о зачислении и передаёт их смотрителю на отправку; после прибытия первокурсников директор вносит решение Распределяющей шляпы, система привязывает ученика к факультету и расписанию, декан заселяет ученика в спальню.}
\ucrow{Главный актор}{Директор}
\ucrow{Второстепенные акторы}{Декан}
\ucrow{Предусловия}{Спальни факультетов заведены в системе.}
\ucrow{Основной поток}{\steps{
    \item Директор формирует список будущих первокурсников и утверждает его.
  \item Система формирует по утверждённому списку письма о зачислении и передаёт их в очередь отправки смотрителя.
  \item После прибытия первокурсников директор открывает список нераспределённых учеников.
  \item Выбирает ученика и вносит факультет, названный Распределяющей шляпой.
  \item Система проверяет, что ученик ещё не распределён, сохраняет факультет и дату распределения.
  \item Система привязывает ученика к расписанию его курса и факультета и к рейтингу факультета.
  \item Декан факультета открывает ученика, система показывает спальни нужного пола и курса со свободными местами.
  \item Декан выбирает спальню, система сохраняет заселение.}}
\ucrow{Альтернативные потоки}{Ученик уже распределён — система отклоняет повторное распределение.\newline
Подходящих спален со свободными местами нет — система сообщает об этом декану, заселение невозможно.\newline
Место в спальне занял другой пользователь одновременно — система отклоняет заселение и возвращает декана к шагу 7.}
\ucrow{Постусловия}{Письма о зачислении переданы на отправку; ученик закреплён за факультетом на всё время обучения и за спальней на текущий учебный год.}
\ucrow{Точки расширения}{Нет свободных мест.}
\end{uctable}
\captionof{table}{Описание прецедента UC-02}
 
\begin{figure}[H]
\centering
\adjustbox{max width=\linewidth}{%
\begin{tikzpicture}
  \node[uc] (l) at (0,7.4)   {Утвердить список\\первокурсников};
  \node[uc] (w) at (9.0,7.4) {Сформировать письма\\о зачислении};
  \node[uc] (a) at (0,4.0)   {Внести решение\\Распределяющей шляпы};
  \node[uc] (b) at (9.0,5.2) {Проверить, что ученик\\не распределён};
  \node[uc] (c) at (9.0,3.0) {Привязать ученика\\к расписанию факультета};
  \node[uc] (d) at (0,0.0)   {Заселить ученика\\в спальню};
  \node[uc] (e) at (9.0,0.8) {Подобрать спальни\\со свободными местами};
  \node[uc] (f) at (9.0,-1.4){Сообщить\\об отсутствии мест};
  \node[sys, fit=(l)(w)(f), label={[font=\small\bfseries]above:UC-02. Распределение и заселение первокурсника}] {};
  \actor{dep}{(-6,5.6)}{Директор}
  \actor{dean}{(-6,-0.2)}{Декан}
  \draw[assoc] (dep-r) -- (l.west);
  \draw[assoc] (dep-r) -- (a.west);
  \draw[assoc] (dean-r) -- (d.west);
  \draw[rel] (l) -- node[lbl,above]{«include»} (w);
  \draw[rel] (a) -- node[lbl,above,sloped]{«include»} (b);
  \draw[rel] (a) -- node[lbl,below,sloped]{«include»} (c);
  \draw[rel] (d) -- node[lbl,above,sloped]{«include»} (e);
  \draw[rel] (f) -- node[cond,below,sloped]{«extend»\\{}[нет свободных мест]} (d);
\end{tikzpicture}}
\caption{Диаграмма прецедента UC-02}
\end{figure}
% ---------------------------------------------------------------------
\subsection{Нарушение и отработка}
 
\begin{uctable}
\ucrow{ID}{UC-03}
\ucrow{Название}{Нарушение и отработка}
\ucrow{Краткое описание}{Преподаватель или смотритель регистрирует нарушение и при необходимости назначает отработку, преподаватель снимает очки с факультета нарушителя; смотритель отмечает выполнение отработки.}
\ucrow{Главный актор}{Преподаватель (в том числе декан) или смотритель}
\ucrow{Второстепенные акторы}{Смотритель (принимает отработку), декан (контролирует нарушения своего факультета)}
\ucrow{Предусловия}{Актор авторизован. Ученик распределён на факультет.}
\ucrow{Основной поток}{\steps{
  \item Актор выбирает ученика, тип нарушения, место и время, описывает причину.
  \item Если актор — преподаватель, он указывает, сколько очков снять с факультета нарушителя.
  \item Если нарушение требует отработки, актор назначает её: дата, время, место, вид работы.
  \item Система проверяет права актора и занятость ученика и в одной транзакции сохраняет нарушение, снятие очков и отработку; рейтинг обновляется, отработка появляется в расписании ученика.
  \item В назначенное время смотритель отмечает выполнение отработки.}}
\ucrow{Альтернативные потоки}{Нарушение возникло автоматически (неявка, конфискат, опоздание) — его фиксирует смотритель, данные уже заполнены.\newline
Нарушение фиксирует смотритель — шаг 2 пропускается, очки по нарушению может снять декан факультета нарушителя.\newline
Время отработки пересекается с уроком или другой отработкой — система показывает конфликт, актор выбирает другое время.\newline
Ученик не явился — отработка получает статус «пропущена», система регистрирует новое нарушение, и сценарий повторяется с шага 1.}
\ucrow{Постусловия}{Нарушение и изменение очков сохранены в журнале; отработка назначена и выполнена либо по неявке зарегистрировано новое нарушение.}
\ucrow{Точки расширения}{Снятие очков, назначение отработки, неявка.}
\end{uctable}
\captionof{table}{Описание прецедента UC-03}
 
\begin{figure}[H]
\centering
\adjustbox{max width=\linewidth}{%
\begin{tikzpicture}
  \node[uc] (a) at (0,5.0)    {Зафиксировать\\нарушение};
  \node[uc] (k) at (0,2.4)    {Просмотреть нарушения\\своего факультета};
  \node[uc] (c) at (0,-0.4)   {Отметить выполнение\\отработки};
  \node[uc] (d) at (9.0,6.4)  {Снять очки\\с факультета};
  \node[uc] (b) at (9.0,3.8)  {Назначить\\отработку};
  \node[uc] (f) at (9.0,1.6)  {Проверить\\занятость ученика};
  \node[uc] (g) at (9.0,-0.4) {Зарегистрировать\\повторное нарушение};
  \node[sys, fit=(a)(d)(c)(g), label={[font=\small\bfseries]above:UC-03. Нарушение и отработка}] {};
  \actor{prof}{(-6.5,5.6)}{Преподаватель}
  \actor{dean}{(-6.5,2.4)}{Декан}
  \actor{filch}{(-6.5,-0.8)}{Смотритель}
  \draw[gen] (-6.5,3.1) -- (-6.5,4.6);
  \draw[assoc] (prof-r) -- (a.west);
  \draw[assoc] (dean-r) -- (k.west);
  \draw[assoc] (filch-r) -- (a.south west);
  \draw[assoc] (filch-r) -- (c.west);
  \draw[rel] (d) -- node[cond,above,sloped]{«extend»\\{}[актор — преподаватель]} (a);
  \draw[rel] (b) -- node[cond,below,sloped]{«extend»\\{}[нарушение требует отработки]} (a);
  \draw[rel] (b) -- node[lbl,right]{«include»} (f);
  \draw[rel] (g) -- node[cond,above]{«extend»\\{}[неявка]} (c);
\end{tikzpicture}}
\caption{Диаграмма прецедента UC-03}
\end{figure}
 

\subsection{Совиная почта и досмотр посылок}

\begin{uctable}
\ucrow{ID}{UC-04}
\ucrow{Название}{Совиная почта и досмотр посылок}
\ucrow{Краткое описание}{Смотритель ведёт совятню: регистрирует входящие посылки ученикам и досматривает их, изымая запрещённые предметы, а также отправляет письма от имени школы, для которых система подбирает свободную сову.}
\ucrow{Главный актор}{Смотритель}
\ucrow{Второстепенные акторы}{Нет}
\ucrow{Предусловия}{Совы и справочник запрещённых предметов заведены в системе.}
\ucrow{Основной поток}{\begin{minipage}[t]{\linewidth}
  \textit{Приём посылки:}
  \begin{enumerate}[nosep,leftmargin=*]
  \item Смотритель регистрирует посылку ученику, доставленную совой: отправитель, получатель, вес, сова.
  \item Система переводит сову в статус «отдыхает» и ставит посылку в очередь досмотра.
  \item Смотритель досматривает посылку по справочнику запрещённых предметов.
  \item Смотритель отмечает посылку как проверенную, система переводит её в статус «допущена к выдаче».
  \end{enumerate}
  \textit{Отправка письма от имени школы:}
  \begin{enumerate}[nosep,leftmargin=*,resume]
  \item Смотритель берёт письмо из очереди отправки (например, письмо о зачислении, сформированное директором) или регистрирует новое: получатель, вес.
  \item Система подбирает свободную сову, способную унести письмо, и переводит её в статус «в полёте».
  \item После возвращения совы смотритель отмечает доставку, сова переводится в статус «отдыхает».
  \end{enumerate}\vspace{1mm}\end{minipage}}
\ucrow{Альтернативные потоки}{Найден запрещённый предмет — смотритель изымает его. Система в одной транзакции заносит предмет в реестр конфиската и регистрирует нарушение на получателя; посылка без изъятого предмета допускается к выдаче.\newline
Свободной совы нужной грузоподъёмности нет — письмо ставится в очередь и отправляется, когда сова освободится.\newline
Письмо не доставлено — смотритель отмечает возврат.}
\ucrow{Постусловия}{Посылка досмотрена и допущена к выдаче, изъятые предметы в реестре конфиската, нарушение зарегистрировано; письмо доставлено.}
\ucrow{Точки расширения}{Изъятие запрещённого предмета, нет свободной совы.}
\end{uctable}
\captionof{table}{Описание прецедента UC-04}

\begin{figure}[H]
\centering
\adjustbox{max width=\linewidth}{%
\begin{tikzpicture}
  \node[uc] (a) at (0,5.6)    {Зарегистрировать\\входящую посылку};
  \node[uc] (d) at (9.0,5.6)  {Поставить посылку\\в очередь досмотра};
  \node[uc] (b) at (0,2.6)    {Досмотреть\\посылку};
  \node[uc] (e) at (9.0,3.4)  {Изъять запрещённый\\предмет};
  \node[uc] (f) at (9.0,1.2)  {Зарегистрировать\\нарушение};
  \node[uc] (m) at (0,-1.0)   {Отправить письмо\\от имени школы};
  \node[uc] (s) at (9.0,-1.0) {Подобрать свободную\\сову по весу};
  \node[uc] (q) at (9.0,-3.2) {Поставить письмо\\в очередь};
  \node[sys, fit=(a)(d)(m)(q), label={[font=\small\bfseries]above:UC-04. Совиная почта и досмотр посылок}] {};
  \actor{filch}{(-6,2.4)}{Смотритель}
  \draw[assoc] (filch-r) -- (a.west);
  \draw[assoc] (filch-r) -- (b.west);
  \draw[assoc] (filch-r) -- (m.west);
  \draw[rel] (a) -- node[lbl,above]{«include»} (d);
  \draw[rel] (e) -- node[cond,above,sloped]{«extend»\\{}[найден запрещённый предмет]} (b);
  \draw[rel] (e) -- node[lbl,right]{«include»} (f);
  \draw[rel] (m) -- node[lbl,above]{«include»} (s);
  \draw[rel] (q) -- node[cond,below,sloped]{«extend»\\{}[нет свободной совы]} (m);
\end{tikzpicture}}
\caption{Диаграмма прецедента UC-04}
\end{figure}
 
 
% ---------------------------------------------------------------------
\subsection{Лечение и пополнение запаса зелий}
 
\begin{uctable}
\ucrow{ID}{UC-05}
\ucrow{Название}{Лечение и пополнение запаса зелий}
\ucrow{Краткое описание}{Целитель регистрирует травму, госпитализирует и лечит ученика зельями; при нехватке зелья зельевар варит его по заявке из ингредиентов теплиц или Хогсмида.}
\ucrow{Главный актор}{Целитель}
\ucrow{Второстепенные акторы}{Зельевар}
\ucrow{Предусловия}{Рецепты зелий и минимальные запасы заведены в системе.}
\ucrow{Основной поток}{\steps{
  \item Целитель регистрирует обращение: пациент, травма, причина, лечение.
  \item Целитель госпитализирует ученика, система закрепляет за ним свободную койку и отмечает освобождение от занятий.
  \item Целитель указывает выданные зелья, система списывает их с запаса.
  \item Запас зелья опустился ниже минимума — система предлагает создать заявку, целитель создаёт её.
  \item Зельевар открывает заявку, система показывает рецепт и потребность в ингредиентах.
  \item Зельевар отмечает варку. Система списывает ингредиенты, пополняет запас зелья и закрывает заявку.}}
\ucrow{Альтернативные потоки}{Госпитализация не нужна — шаг 2 пропускается.\newline
Свободных коек нет — система сообщает об этом, госпитализация невозможна.\newline
Запас не ниже минимума — сценарий завершается после шага 3.\newline
Ингредиентов в запасе не хватает — заявка получает статус «ожидает ингредиентов». Зельевар вносит урожай из теплиц или, если ингредиент в теплицах не растёт, заказывает его в Хогсмиде. После поступления сценарий продолжается с шага 6.}
\ucrow{Постусловия}{Обращение в медицинском журнале, освобождение видно преподавателям, запас зелья восстановлен.}
\ucrow{Точки расширения}{Госпитализация, заявка на варку, нехватка ингредиентов.}
\end{uctable}
\captionof{table}{Описание прецедента UC-05}
 
\begin{figure}[H]
\centering
\adjustbox{max width=\linewidth}{%
\begin{tikzpicture}
  \node[uc] (a) at (0,5.2)    {Зарегистрировать\\обращение};
  \node[uc] (d) at (9.0,7.0)  {Госпитализировать\\на свободную койку};
  \node[uc] (e) at (9.0,4.8)  {Списать выданные\\зелья};
  \node[uc] (f) at (9.0,2.6)  {Создать заявку\\на варку};
  \node[uc] (b) at (0,0.4)    {Сварить зелье\\по заявке};
  \node[uc] (i) at (9.0,0.4)  {Проверить запас\\ингредиентов};
  \node[uc] (g) at (9.0,-1.8) {Заказать ингредиенты\\в Хогсмиде};
  \node[uc] (c) at (0,-2.8)   {Зарегистрировать урожай\\из теплиц};
  \node[sys, fit=(a)(d)(c)(g), label={[font=\small\bfseries]above:UC-05. Лечение и пополнение запаса зелий}] {};
  \actor{heal}{(-6,5.0)}{Целитель}
  \actor{pot}{(-6,-1.2)}{Зельевар}
  \draw[assoc] (heal-r) -- (a.west);
  \draw[assoc] (pot-r) -- (b.west);
  \draw[assoc] (pot-r) -- (c.west);
  \draw[rel] (d) -- node[cond,above,sloped]{«extend»\\{}[нужна госпитализация]} (a);
  \draw[rel] (a) -- node[lbl,above]{«include»} (e);
  \draw[rel] (f) -- node[cond,below,sloped]{«extend»\\{}[запас ниже минимума]} (a);
  \draw[rel] (b) -- node[lbl,above]{«include»} (i);
  \draw[rel] (g) -- node[cond,below,sloped]{«extend»\\{}[ингредиента нет в теплицах]} (b);
\end{tikzpicture}}
\caption{Диаграмма прецедента UC-05}
\end{figure}
% ---------------------------------------------------------------------
\subsection{Допуск и выход в Хогсмид}
 
\begin{uctable}
\ucrow{ID}{UC-06}
\ucrow{Название}{Допуск и выход в Хогсмид}
\ucrow{Краткое описание}{Декан утверждает разрешение опекуна, система проверяет допуск и формирует список, смотритель отмечает выход и возвращение учеников.}
\ucrow{Главный актор}{Декан}
\ucrow{Второстепенные акторы}{Смотритель}
\ucrow{Предусловия}{Директор объявил выход в Хогсмид. Разрешение опекуна пришло совиной почтой.}
\ucrow{Основной поток}{\steps{
  \item Декан открывает список учеников своего факультета 3 курса и старше.
  \item Декан регистрирует полученное разрешение и утверждает его.
  \item Система проверяет допуск: курс, утверждённое разрешение, отсутствие невыполненных отработок.
  \item Система включает ученика в список допущенных.
  \item В день выхода смотритель проверяет ученика по списку и отмечает выход.
  \item При возвращении смотритель отмечает время возвращения.}}
\ucrow{Альтернативные потоки}{У ученика есть невыполненная отработка или декан запретил выход — ученик не допускается, причина сохраняется и видна декану.\newline
Ученика нет в списке — система не допускает отметку выхода.\newline
Ученик вернулся позже назначенного времени — система регистрирует нарушение.}
\ucrow{Постусловия}{Известно, кто вышел и вернулся; опоздания переданы в дисциплинарный цикл.}
\ucrow{Точки расширения}{Запрет выхода, опоздание.}
\end{uctable}
\captionof{table}{Описание прецедента UC-06}
 
\begin{figure}[H]
\centering
\adjustbox{max width=\linewidth}{%
\begin{tikzpicture}
  \node[uc] (a) at (0,4.2)    {Утвердить разрешение\\опекуна};
  \node[uc] (c) at (9.0,5.4)  {Проверить допуск\\ученика};
  \node[uc] (d) at (9.0,2.9)  {Запретить выход};
  \node[uc] (b) at (0,0.2)    {Отметить выход\\и возвращение};
  \node[uc] (e) at (9.0,1.0)  {Проверить ученика\\по списку допущенных};
  \node[uc] (f) at (9.0,-1.2) {Зарегистрировать\\нарушение};
  \node[sys, fit=(a)(c)(f), label={[font=\small\bfseries]above:UC-06. Допуск и выход в Хогсмид}] {};
  \actor{dean}{(-6,4.0)}{Декан}
  \actor{filch}{(-6,0.0)}{Смотритель}
  \draw[assoc] (dean-r) -- (a.west);
  \draw[assoc] (filch-r) -- (b.west);
  \draw[rel] (a) -- node[lbl,above,sloped]{«include»} (c);
  \draw[rel] (d) -- node[cond,below,sloped,pos=0.5]{«extend»\\{}[есть отработка\\или запрет декана]} (a);
  \draw[rel] (b) -- node[lbl,above,sloped]{«include»} (e);
  \draw[rel] (f) -- node[cond,below,sloped]{«extend»\\{}[опоздание]} (b);
\end{tikzpicture}}
\caption{Диаграмма прецедента UC-06}
\end{figure}
 
 
\section{Архитектура системы}
 
\subsection{Общая структура}
 
Система строится по клиент-серверной архитектуре. Пользователь работает в браузере,
браузер отправляет запросы на сервер, сервер обрабатывает их, обращается к базе данных
и возвращает готовую HTML-страницу. Страницы формируются на сервере шаблонизатором
Thymeleaf, поэтому отдельного клиентского приложения нет.
 
\begin{figure}[H]
\centering
\adjustbox{max width=\linewidth}{%
\begin{tikzpicture}[
  comp/.style={draw, fill=black!4, rounded corners=1pt, align=center, font=\small,
               minimum height=11mm, text width=22mm, inner sep=2pt},
  grp/.style={draw, rounded corners=2pt, inner sep=3mm},
  glbl/.style={font=\small\bfseries},
  arr/.style={-{Stealth[length=2mm]}},
  lab/.style={font=\scriptsize, fill=white, inner sep=1pt}
]
  % браузер
  \node[comp, text width=22mm] (br) at (-0.6,0) {HTML-страницы\\и формы};
  \node[grp, fit=(br), label={[glbl]above:Браузер}] (brg) {};
 
  % веб-приложение
  \node[comp] (sec) at (3.9,0)  {Spring Security};
  \node[comp] (ctl) at (7.0,0.9) {Controllers\\(Spring MVC)};
  \node[comp] (tpl) at (7.0,-0.9) {Шаблоны\\Thymeleaf};
  \node[grp, fit=(ctl)(tpl), label={[glbl]above:Presentation Layer}] (pres) {};
 
  \node[comp] (svc) at (10.6,0)  {Business Services\\(\texttt{@Service},\\ \texttt{@Transactional})};
  \node[grp, fit=(svc), label={[glbl]above:Application Layer}] (app) {};
 
  \node[comp] (rep) at (14.0,0)  {Spring Data JPA\\Repositories};
  \node[comp] (jpa) at (16.9,0)  {JPA /\\Hibernate};
  \node[grp, fit=(rep)(jpa), label={[glbl]above:Persistence Layer}] (pers) {};
 
  \node[grp, fit=(sec)(pres)(app)(pers), inner sep=6mm, label={[glbl]above:Web Application}] (web) {};
  \node[grp, fit=(web), inner sep=6mm, label={[glbl]above:Spring Boot}] (boot) {};
  \node[grp, fit=(boot), inner sep=6mm, label={[glbl]above:helios}] (hel) {};
 
  % БД
  \node[cylinder, shape border rotate=90, draw, fill=black!4, aspect=0.25,
        minimum height=14mm, minimum width=20mm, align=center, font=\small]
        (db) at (22.4,0) {PostgreSQL};
 
  \draw[arr, <->] (brg.east) -- node[lab,above]{HTTP / HTML} node[lab,below]{SSH-туннель} (sec.west);
  \draw[arr] (sec.east) -- (pres.west);
  \draw[arr] (ctl) -- (tpl);
  \draw[arr] (pres.east) -- (app.west);
  \draw[arr] (app.east) -- (pers.west);
  \draw[arr] (rep) -- (jpa);
  \draw[arr, <->] (jpa.east) -- node[lab,above]{JDBC} (db.west);
\end{tikzpicture}}
\caption{Общая структура системы}
\end{figure}
 
Запрос из браузера проходит через фильтры Spring Security (проверка сеанса и роли)
и попадает в DispatcherServlet, который передаёт его нужному контроллеру. Контроллер
вызывает сервис бизнес-логики, кладёт результат в модель и возвращает имя
представления; ViewResolver находит шаблон Thymeleaf, который формирует HTML-страницу.
Сервисы работают с данными через репозитории Spring Data JPA, а Hibernate переводит
операции с сущностями в SQL-запросы к PostgreSQL.
 
Серверная часть — монолитное приложение, разделённое на модули по бизнес-процессам.
Все модули работают в составе одного приложения и используют общую базу данных.
\begin{itemize}
  \item \textbf{Модуль пользователей и доступа} — вход и выход, роли, учётные записи,
        проверка прав доступа к данным своего факультета.
  \item \textbf{Модуль учеников и расписания} — ученики, сотрудники, спальни,
        расписание занятий; предоставляет эти данные остальным модулям.
  \item \textbf{Модуль распределения} — решение Распределяющей шляпы,
        подбор спальни и заселение.
  \item \textbf{Модуль дисциплины} — нарушения, очки факультетов и их журнал,
        отработки, рейтинг и Кубок школы. Его сервис вызывают модули почты и Хогсмида,
        когда нужно зарегистрировать нарушение.
  \item \textbf{Модуль совиной почты} — совы, отправления и рассылки, очередь досмотра,
        реестр конфиската.
  \item \textbf{Модуль больничного крыла} — обращения, койки, запас зелий,
        заявки на варку, рецепты, ингредиенты и заказы поставщикам.
  \item \textbf{Модуль Хогсмида} — выходы, разрешения опекунов, проверка допуска,
        отметки выхода и возвращения.
  \item \textbf{Модуль отчётов} — отчёты для директора и деканов.
\end{itemize}
 
\subsection{Используемые технологии}
 
\begin{longtable}{|>{\raggedright\arraybackslash}p{34mm}|>{\raggedright\arraybackslash}p{48mm}|>{\raggedright\arraybackslash}p{72mm}|}
\hline
\rowcolor{black!8}\textbf{Слой} & \textbf{Технология} & \textbf{Обоснование} \\ \hline
\endhead
Frontend & \textbf{Spring MVC + Thymeleaf}, CSS & Страницы формируются на сервере, отдельный фронтенд-проект и Node.js не нужны \\ \hline
Backend & \textbf{Spring Boot 3}  & Приложение на Spring Boot само содержит контейнер. Поэтому достаточно запустить один jar \\ \hline
Бизнес-логика & Бины-сервисы Spring (\texttt{@Service}) & Каждый бизнес-процесс — отдельный модуль; контроллер вызывает сервис, сервис вызывает репозиторий \\ \hline
Доступ к данным & Spring Data JPA (\texttt{@Repository}) + Hibernate; нативные запросы и вызовы функций PL/pgSQL; пул соединений HikariCP & Типовые операции — методами репозиториев, сложные — запросами JPQL и SQL \\ \hline
БД & PostgreSQL & Требование задания \\ \hline
Сборка & Gradle с плагином Spring Boot & Задача \texttt{bootJar} упаковывает приложение, шаблоны и статику в один \texttt{jar} \\ \hline
Язык & Java 17 & Минимальная версия для Spring Boot 3 \\ \hline
\end{longtable}
 
\end{document}